\documentclass[aspectratio=169,10pt,spanish]{beamer}

\input{../../../_preambulo_beamer.tex}
\input{../../../_comandos_beamer.tex}

\selectlanguage{spanish}

\title{MA1001B - Modelación Estadística para la Toma de Decisiones}
\subtitle{Lección 28: Intervalo t de Student cuando sigma es desconocida}
\author{}
\institute{Tecnol\'ogico de Monterrey}
\date{}

\begin{document}

\begin{frame}[plain]
  \vspace{1.2cm}
  {\Large\bfseries MA1001B - Modelación Estadística para la Toma de Decisiones\par}
  \vspace{0.7cm}
  {\large Lección 28: Intervalo t de Student cuando sigma es desconocida\par}
  \vspace{0.7cm}
  {\color{TecRojo}\rule{\textwidth}{1.2pt}}
  \vspace{0.8cm}

  Facultad MA1001B\\[0.3cm]
  Periodo\\[0.3cm]
  Tecnol\'ogico de Monterrey
\end{frame}

\begin{frame}{La pregunta del intervalo con sigma desconocida}
  \begin{alertblock}{Pregunta guía}
    ¿Cómo se construye un intervalo al 95\% para una media desconocida
    cuando $\sigma$ debe sustituirse por la $s$ muestral?
  \end{alertblock}

  \vspace{0.6em}
  Al final de esta lección, usted puede:
  \begin{itemize}
    \item calcular $s$ y $\widehat{\mathrm{EE}}=s/\sqrt{n}$;
    \item obtener $t^{*}$ de la distribución $t$ con $\mathrm{gl}=n-1$;
    \item formar $\bar{x}\pm t^{*}\cdot s/\sqrt{n}$;
    \item explicar cómo un atípico infla $s$ y el intervalo.
  \end{itemize}

  \vfill
  \scriptsize Todos los volúmenes de llenado usados hoy son sintéticos. No se usan datos reales de empresa.
\end{frame}

\begin{frame}{La misma muestra, ahora con $\sigma$ desconocida}
  \small
  \begin{center}
    \begin{tabular}{lr}
      \toprule
      \textbf{Cantidad} & \textbf{Valor} \\
      \midrule
      Tamaño de muestra $n$ & 36 botellas \\
      Grados de libertad $n-1$ & 35 \\
      Estimación puntual $\bar{x}$ (semilla 42) & 502.738498 ml \\
      $s$ muestral & 8.390375 ml \\
      \bottomrule
    \end{tabular}
  \end{center}

  \vspace{0.5em}
  \begin{exampleblock}{Error estándar estimado}
    \[
      \widehat{\mathrm{EE}}=\frac{s}{\sqrt{n}}=\frac{8.390375}{\sqrt{36}}=1.398396.
    \]
    La incertidumbre extra en $s$ es la razón de que el valor crítico sea
    $t$, no $z$.
  \end{exampleblock}
\end{frame}

\begin{frame}[fragile]{Paso 1: Desviación estándar muestral}
  \begin{columns}[T]
    \begin{column}{0.42\textwidth}
      \small
      \begin{lstlisting}
s = np.std(sample, ddof=1)
se_hat = s / np.sqrt(n)
df = n - 1
      \end{lstlisting}

      \begin{block}{Salida verificada}
        $\mathrm{gl}=35$\\
        $s=8.390375$ ml\\
        $\widehat{\mathrm{EE}}=1.398396$
      \end{block}
    \end{column}
    \begin{column}{0.58\textwidth}
      \centering
      \includegraphics[width=\textwidth]{../figuras/t_ci_mean_01_sample_sd.png}
      \vspace{0.2em}
      \scriptsize Histograma y banda de EE estimado del Paso 1
    \end{column}
  \end{columns}
\end{frame}

\begin{frame}{Sustituir $z^{*}$ por $t^{*}$ con $\mathrm{gl}=35$}
  \small
  De \texttt{scipy.stats.t.ppf}:
  \[
    t^{*}=t_{0.975,35}=2.030108 > z^{*}=1.96.
  \]
  Margen de error e intervalo:
  \[
    E=2.030108\times 1.398396=2.838894,
  \]
  \[
    \bar{x}\pm E=[499.899604, 505.577392].
  \]
  Este intervalo con semilla 42 cubre la $\mu$ oculta $=502$. El valor
  crítico $t$ es mayor que $1.96$; $s$ misma todavía puede ser menor que
  $\sigma$.
\end{frame}

\begin{frame}[fragile]{Paso 2: El intervalo t al 95\%}
  \begin{columns}[T]
    \begin{column}{0.40\textwidth}
      \small
      \begin{lstlisting}
t_star = stats.t.ppf(
    0.975, df=35
)
margin = t_star * se_hat
      \end{lstlisting}

      \begin{block}{Salida verificada}
        $t^{*}=2.030108$\\
        Margen $=2.838894$\\
        IC $[499.899604, 505.577392]$
      \end{block}
    \end{column}
    \begin{column}{0.60\textwidth}
      \centering
      \includegraphics[width=\textwidth]{../figuras/t_ci_mean_02_t_interval.png}
      \vspace{0.2em}
      \scriptsize Intervalo t con sigma desconocida del Paso 2
    \end{column}
  \end{columns}
\end{frame}

\begin{frame}{Paso 3: Un atípico infla $s$ y el intervalo}
  \begin{columns}[T]
    \begin{column}{0.42\textwidth}
      \small
      Un llenado sustituido por $560$ ml.

      \vspace{0.4em}
      $s$ sube de $8.390375$ a $12.616189$.

      \vspace{0.4em}
      Intervalo contaminado:
      \[
        [499.767301, 508.304710].
      \]
      La anchura crece $2.859620$ ml.

      \vspace{0.5em}
      \textbf{Límite:} una observación extrema infla $s$ y el intervalo $t$.
      La fórmula sigue corriendo; la estimación de escala no.
    \end{column}
    \begin{column}{0.58\textwidth}
      \centering
      \includegraphics[width=\textwidth]{../figuras/t_ci_mean_03_outlier_inflates_s.png}
      \vspace{0.2em}
      \scriptsize Intervalos t limpio versus inflado por atípico del Paso 3
    \end{column}
  \end{columns}
\end{frame}

\begin{frame}{Verificación de conceptos}
  \begin{enumerate}
    \item ¿Por qué el valor crítico es $t_{0.975,35}$ en lugar de $1.96$?
    \item Calcule $\widehat{\mathrm{EE}}=8.390375/\sqrt{36}$.
    \item ¿Por qué este intervalo $t$ todavía puede ser más estrecho que el
      intervalo z de la Lección 27?
    \item ¿Qué le ocurre a $s$ cuando un llenado se sustituye por $560$ ml?
  \end{enumerate}

  \vspace{0.6em}
  \begin{block}{Conexión con la Lección 29}
    La sesión puede continuar directamente con un intervalo de confianza
    para una proporción poblacional y las condiciones $np$, $n(1-p)$.
  \end{block}

  \vfill
  \scriptsize Fuente: OpenStax, \textit{Introductory Business Statistics 2e},
  Capítulo 8, Sección 8.2. CC BY-NC-SA.
\end{frame}

\appendix



\end{document}
